%INTRO


\begin{langblock}{0.425}
	\begin{eng}
	
	Riding a motorcycle is a demanding and complex task. Skilled riders pay attention to the riding environment and to operating the motorcycle, identifying potential hazards, making good judgments, and executing decisions quickly and skillfully. Your ability to perform and respond to changing road and traffic conditions is influenced by how fit and alert you are. Alcohol and other drugs, more than any other factor, degrade your ability to think clearly and to ride safely.  As little as one drink can have a significant effect on your performance.
	
	Let's look at the risks involved in riding after drinking or using drugs. What to do to protect yourself and your fellow riders is also examined.
	\end{eng}
	\begin{rus}
		
		Управление мотоциклом "--- это сложная и ответственная задача. Опытные мотоциклисты уделяют максимум внимания обстановке на дороге и управлению мотоциклом, выявляя потенциальные опасности, принимая взвешенные решения и быстро и умело их выполняя. На вашу способность действовать и реагировать на меняющиеся дорожные и дорожно-транспортные условия влияет ваша физическая форма и бдительность. Алкоголь и другие наркотические вещества, больше чем какой либо другой фактор, ухудшают вашу способность ясно мыслить и безопасно управлять мотоциклом. Даже одна порция алкоголя может существенно повлиять на ваши способности.
		
		Давайте рассмотрим риски, связанные с вождением после употребления алкоголя или наркотиков. Также будет рассмотрено, что нужно делать, чтобы защитить себя и других мотоциклистов.
\end{rus}
\end{langblock}




\langsection[subsection]{Why This Information Is Important}{Почему это важно}


\begin{langblock}{0.425}
	\begin{eng}
	
	Alcohol is a major contributor to motorcycle crashes, particularly fatal crashes. Studies show that 40\% to 45\% of all riders killed in motorcycle crashes had been drinking. Only one-third of those riders had a blood alcohol concentration above legal limits. The rest had only a few drinks in their  systems --- enough to impair riding skills. In the past, drug levels have been harder to distinguish or have not been separated from drinking violations for the traffic records. But riding "under the influence" of either alcohol or drugs poses physical and legal hazards for every rider.
	
	Drinking and drug use is as big a problem among motorcyclists as it is among automobile drivers.  Motorcyclists, however, are more likely to be killed or severely injured in a crash. Injuries occur in 90\% of motorcycle crashes and 33\% of automobile crashes that involve abuse of substances. On a yearly basis, 2,100 motorcyclists are killed and about 50,000 seriously injured in this same type of crash. These statistics are too overwhelming to ignore.
	
	
	By becoming knowledgeable about the effects of alcohol and other drugs you will see that riding and substance abuse don't mix. Take positive steps to protect yourself and prevent others from injuring themselves.
	\end{eng}
	\begin{rus}
		
		Алкоголь является одной из основных причин ДТП с участием мотоциклов, особенно с летальным исходом. Исследования показывают, что от 40\% до 45\% всех мотоциклистов, погибших в ДТП, употребляли алкоголь. Только у трети из них концентрация алкоголя в крови превышала установленные законом пределы. У остальных в организме было всего несколько порций алкоголя "--- но этого достаточно, чтобы ухудшить способности к вождению. Раньше уровень наркотиков в организме было сложнее определить, либо он не выделялся отдельно от нарушений, связанных с употреблением алкоголя, в данных о дорожно-транспортных происшествиях. Однако вождение под алкогольным или наркотическим опьянением представляет физическую и юридическую опасность для каждого мотоциклиста.
		
		Употребление алкоголя и наркотиков является столь же серьезной проблемой среди мотоциклистов, как и среди автомобилистов. Однако мотоциклисты с большей вероятностью погибают или получают тяжелые травмы в результате ДТП. Травмы фиксируются в 90\% ДТП с участием мотоциклов и в 33\% ДТП с участием автомобилей, связанных со употреблением психоактивных веществ. Ежегодно в результате подобных ДТП гибнут 2'100 мотоциклистов, а около 50'000 получают тяжелые травмы. Эти статистические данные слишком впечатляющи, чтобы их игнорировать.
		
		
		Узнав о воздействии алкоголя и других наркотических веществ, вы поймете, что вождение мотоцикла и злоупотребление психоактивными веществами несовместимы. Примите активные меры, чтобы защитить себя и предотвратить травмы других людей.
\end{rus}
\end{langblock}


\langsection[subsection]{Alcohol and Other Drugs in Motorcycle Operation}{Алкоголь и прочие препараты}



\begin{langblock}{0.425}
	\begin{eng}
	
	No one is immune to the effects of alcohol or drugs. Friends may brag about their ability to hold their liquor or perform better on drugs, but alcohol or drugs make them less able to think clearly and perform physical tasks skillfully. Judgment and the decisionmaking processes needed for vehicle operation are affected long before legal limitations are reached.
	
	
	Many over-the-counter, prescription, and illegal drugs have side effects that increase the risk of riding. It is difficult to accurately measure the involvement of particular drugs in motorcycle crashes. But we do know what effects various drugs have on the process involved in riding a motorcycle. We also know that the combined effects of alcohol and other drugs are more dangerous than either is alone.
	\end{eng}
	\begin{rus}
		
		Не бывает стойких к воздействию алкоголя или наркотиков людей. Друзья могут хвастаться своей выносливостью к алкоголю или повышенной работоспособностью под воздействием наркотиков, но алкоголь и наркотики снижают их способность ясно мыслить и выполнять физические задачи. На способность адекватно оценивать ситуацию и принимать решения, необходимые для управления транспортным средством, это влияет задолго до того, как будут достигнуты установленные законом пределы.
		
		
		Многие безрецептурные, рецептурные и наркотические препараты имеют побочные эффекты, которые повышают риск при вождении мотоцикла. Точно определить степень влияния конкретных препаратов на аварии с участием мотоциклов довольно сложно. Однако мы знаем, какое воздействие различные наркотики оказывают на процессы, связанные с управлением мотоциклом. Мы также знаем, что сочетание воздействия алкоголя и других веществ более опасно, чем воздействие каждого из них в отдельности.
\end{rus}
\end{langblock}





\langsection[subsection]{Alcohol in the Body}{Алкоголь в организме}



\begin{langblock}{0.425}
	\begin{eng}
	
	Alcohol enters the bloodstream quickly. Unlike most foods and beverages, it does not need to be digested. Within minutes after being consumed, it reaches the brain and begins to affect the drinker. The major effect alcohol has is to slow down and impair bodily functions --- both mental and physical. Whatever you do, you do less well after consuming alcohol.
	\end{eng}
	\begin{rus}
		
		Алкоголь быстро попадает в кровоток. В отличие от большинства продуктов и напитков, он не проходит через пищеварение. Уже через несколько минут после употребления он достигает мозга и начинает оказывать воздействие на выпившего. Основное действие алкоголя заключается в замедлении и ухудшении функций организма "--- как умственных, так и физических. Что бы вы ни делали, после употребления алкоголя вы делаете это хуже.
\end{rus}
\end{langblock}


\langsection[subsubsection]{Blood Alcohol Concentration}{Концентрация алкоголя в крови (BAC)}


\begin{langblock}{0.425}
	\begin{eng}
	
	Blood Alcohol Concentration or BAC is the amount of alcohol in relation to blood in the body.  Generally, alcohol can be eliminated in the body at the rate of almost one drink per hour. But a variety of other factors may also influence the level of alcohol retained. The more alcohol in your blood, the greater the degree of impairment.
	
	Three factors play a major part in determining BAC:
	
	
	\begin{itemize}[noitemsep]
		\item The amount of alcohol you consume.
		\item How fast you drink.
		\item Your body weight.
	\end{itemize}
	
	Other factors also contribute to the way alcohol affects your system. Your sex, physical condition and food intake are just a few that may cause your BAC level to be even higher. But the full effects of these are not completely known. Alcohol may still accumulate in your body even if you are drinking at a rate of one drink per hour. Abilities and judgment can be affected by that one drink.
	
	A 12-ounce can of beer, a mixed drink with one shot of liquor, and a 5-ounce glass of wine all contain the same amount of alcohol.
	
	The faster you drink, the more alcohol accumulates in your body. If you drink two drinks in an hour, at the end of that hour, at least one drink will remain in your bloodstream.
	
	Without taking into account any of the other factors, the formula below illustrates the LEAST amount of drinks remaining in the bloodstream:
	
	Total drinks consumed MINUS number of hours since last drink EQUALS drinks left in body
	
	
	
	A person drinking:
	\begin{itemize}[noitemsep]
		\item 8 drinks in 4 hours would have at least 4 drinks remaining in their system.
		\item 7 drinks in 3 hours would have at least 4 drinks remaining in their system.
	\end{itemize}
	
	There are times when a larger person may not accumulate as high a concentration of alcohol for each drink consumed. They have more blood and other bodily fluids. But because of individual differences it is better not to take the chance that abilities and judgment have not been affected. Whether or not you are legally intoxicated is not the real issue. Impairment of judgment and skills begins well below the legal limit.
	\end{eng}
	\begin{rus}
		
		Концентрация алкоголя в крови (BAC) "--- это количество алкоголя в крови человека. Как правило, алкоголь выводится из организма со скоростью примерно одна порция в час. Однако на уровень алкоголя в крови могут влиять и другие факторы. Чем больше алкоголя в крови, тем сильнее нарушение координации движений.
		
		Три фактора играют основную роль в определении BAC:
		
		
		\begin{itemize}[noitemsep]
			\item Количество употребленного алкоголя.
			\item Скорость употребления алкоголя.
			\item Ваша масса тела.
		\end{itemize}
		
		Другие факторы также влияют на то, как алкоголь воздействует на ваш организм. Ваш пол, физическое состояние и прием пищи "--- это лишь некоторые из факторов, которые могут привести к еще более высокому уровню BAC. Однако полное влияние этих факторов до конца не изучено. Алкоголь может накапливаться в организме даже при употреблении по одной порции в час. Уже один напиток может повлиять на ваши возможности и способность адекватно оценивать ситуацию.
		
		12oz банка пива, коктейль с одной порцией крепкого алкоголя и 5oz бокал вина содержат одинаковое количество алкоголя.
		
		Чем быстрее вы пьете, тем больше алкоголя накапливается в организме. Если вы выпили две порции за час, то к концу этого часа в крови останется как минимум одна порция.
		
		Не учитывая никаких других факторов, приведенная ниже формула показывает МИНИМАЛЬНОЕ количество напитков, остающихся в кровотоке:
		
		Общее количество выпитых напитков МИНУС количество часов с момента последнего напитка РАВНО количество напитков, оставшихся в организме
		
		
		
		Человек, выпивший:
		\begin{itemize}[noitemsep]
			\item 8 порций за 4 часа, будет иметь в организме не менее 4 порций.
			\item 7 порций за 3 часа, будет иметь в организме не менее 4 порций.
		\end{itemize}
		
		Бывают случаи, когда у человека крупного телосложения концентрация алкоголя на каждую выпитую порцию может быть не столь высокой. У таких людей больше крови и других биологических жидкостей. Однако из-за индивидуальных особенностей лучше не рисковать, полагая, что физические возможности и способность к адекватному суждению не пострадали. Находитесь ли вы в состоянии алкогольного опьянения с юридической точки зрения "--- не самая важная проблема. Нарушение способности к адекватному суждению и координации начинается значительно ниже установленного законом предела.
\end{rus}
\end{langblock}


\langsection[subsection]{Penalties for Driving While under the Influence of Alcohol or Drugs}{Штрафы за вождение в нетрезвом виде}



\begin{langblock}{0.425}
	\begin{eng}
	
	The penalties for driving while under the influence of alcohol or drugs are severe. In Pennsylvania you are considered to be driving while under the influence if your blood alcohol level is .08\% or higher. For a first conviction or acceptance into an Accelerated Rehabilitation Disposition (ARD) Program, you could:
	
	\begin{itemize}[noitemsep]
		\item Lose your license.
		\item Be sentenced to jail.
		\item Pay fines.
		\item Be ordered to undergo alcohol and drug rehabilitation treatment and/or Alcohol Highway Safety School.
	\end{itemize}
	\end{eng}
	\begin{rus}
		
		Наказания за вождение в состоянии алкогольного или наркотического опьянения являются суровыми. В Пенсильвании считается, что вы управляете транспортным средством «в состоянии опьянения», если уровень алкоголя в крови составляет 0,08\% или выше. В случае первого осуждения или включения в программу ускоренной реабилитации (ARD) вам грозит:
		
		\begin{itemize}[noitemsep]
			\item Лишение водительских прав.
			\item быть приговоренным к тюремному заключению;
			\item уплатить штрафы;
			\item получить предписание пройти курс реабилитации от алкогольной и наркотической зависимости и/или курсы по безопасности дорожного движения для водителей, употребляющих алкоголь.
		\end{itemize}
\end{rus}
\end{langblock}

\langsection[subsubsection]{Pennsylvania'S Implied Consent Law}{Закон о подразумеваемом согласии}




\begin{langblock}{0.425}
	\begin{eng}
	
	The Implied Consent law is very important to you as a rider.  If the police arrest you for driving while under the influence of alcohol or drugs and you refuse to take one or more chemical tests of breath, blood, or urine, your driving privilege will be automatically suspended for one (1) year in addition to the driving privilege suspension imposed for a conviction or ARD for driving while under the influence. Altogether, a conviction plus refusal could result in a two (2) year driving privilege suspension.
	
	The law covering chemical testing says that you have agreed to take such a test just by being licensed to drive in Pennsylvania.  Even if you are not found guilty of driving while under the influence, your driving privilege will be suspended for one (1) year if you refuse to take a blood, breath, or urine test.
	\end{eng}
	\begin{rus}
		
		Закон «Подразумеваемое согласие» имеет для вас, как водителя, огромное значение. Если полиция задержит вас за вождение в состоянии алкогольного или наркотического опьянения, и вы откажетесь пройти один или несколько тестов на алкоголь в выдыхаемом воздухе, крови или моче, ваше право на вождение будет автоматически приостановлено на один (1) год в дополнение к сроку приостановки, наложенному в результате осуждения или участия в программе ARD за вождение в состоянии опьянения. В совокупности осуждение и отказ от прохождения теста могут привести к приостановке водительских прав на срок два (2) года.
		
		Закон, регулирующий проведение химических тестов, гласит, что вы даете согласие на прохождение такого теста уже тем, что имеете водительские права в штате Пенсильвания.  Даже если вас не признают виновным в вождении в состоянии опьянения, ваши водительские права будут приостановлены на один (1) год, если вы откажетесь пройти тест.
\end{rus}
\end{langblock}

\langsection[subsection]{Minimize the Risks}{Снижение рисков}




\begin{langblock}{0.425}
	\begin{eng}
	
	Your ability to judge how well you are riding is affected first.  Although you may be performing more and more poorly, you think you are doing better and better. The result is that you ride confidently, taking greater and greater risks.  Minimize the risks of drinking and riding by taking steps before you drink. Control your drinking or control your riding.
	\end{eng}
	\begin{rus}
		
		В первую очередь это сказывается на вашей способности правильно оценивать, насколько хорошо вы управляете мотоциклом. Хотя ваши навыки могут ухудшаться, вам кажется, что вы ездите все лучше и лучше. В результате вы чувствуете себя уверенно за рулем и идете на все больший риск. Сведите к минимуму риски, связанные с вождением в нетрезвом виде, приняв меры еще до того, как начнете пить. Сдержите желание выпить, либо откажитесь от поездки.
\end{rus}
\end{langblock}

\clearpage
\langsection[subsection]{Make an Intelligent Choice}{Сделайте правильный выбор}


\begin{langblock}{0.425}
	\begin{eng}
	
	Don't Drink -- Once you start, your resistance becomes weaker
	
	Setting a limit or pacing yourself are poor alternatives at best. Your ability to exercise good judgment is one of the first things affected by alcohol. Even if you have tried to drink in moderation, you may not realize to what extent your skills have suffered from alcohol's fatiguing effects.
	
	
	Or, Don't Ride -- If you haven't controlled your drinking, you must control your riding.
	
	\begin{itemize}[noitemsep]
		\item Leave the motorcycle home -- so you won't be tempted to ride.  Arrange another way to get home.
		\item Wait -- If you exceed your limit, wait until your system eliminates the alcohol and its fatiguing effects.
	\end{itemize}
	\end{eng}
	\begin{rus}
		
	
		Не пейте "--- как только вы начнете, ваша самоконтроль ослабеет
		
		Установка лимита или попытки контролировать темп "--- в лучшем случае неэффективные меры. Способность трезво оценивать ситуацию "--- одна из первых вещей, на которую влияет алкоголь. Даже если вы пытались пить в меру, вы можете не осознавать, насколько ваши навыки  вождения пострадали от воздействия алкоголя.
		
		
		Или если выпили: не садитесь за руль "--- если вы не смогли воздержаться от спиртного, вы должны воздержаться от вождения.
		
		\begin{itemize}[noitemsep]
			\item Оставьте мотоцикл дома "--- так у вас не возникнет соблазна сесть за руль.  Организуйте другой способ добраться до дома.
			\item Подождите "--- если вы превысили допустимый уровень, подождите, пока ваш организм переработает алкоголь и его пагубное воздействие.
		\end{itemize}
\end{rus}
\end{langblock}

\langsection[subsection]{Steps to Protect Friends}{Чтобы защитить друзей}

\begin{langblock}{0.425}
	\begin{eng}
	
	People who have had too much to drink are unable to make a responsible decision. It is up to others to step in and keep them from taking too great a risk. No one wants to do this -- it's uncomfortable, embarrassing and thankless. You are rarely thanked for your efforts at the time. But the alternatives are often worse.
	
	There are several ways to keep friends from hurting themselves:
	
	\begin{itemize}[noitemsep]
		\item Arrange a safe ride -- Provide alternative ways for them to get home.
		\item Slow the pace of drinking -- Involve them in other activities.
		\item Keep them there -- Use any excuse to keep them from getting on their motorcycle. Serve them food and coffee to pass the time.  Explain your concerns for their risks of getting arrested or hurt, or hurting someone else.
		\item Get friends involved -- Use peer pressure from a group of friends to intervene.
	\end{itemize}
	It helps to enlist support from others when you decide to step in. The more people on your side, the easier it is to be firm and the harder it is for the rider to resist. While you may not be thanked at the time, you will never have to say, If only I had.
	\end{eng}
	\begin{rus}
		
		Люди, которые перебрали с алкоголем, неспособны принимать взвешенные решения. И тут окружающие должны вмешаться и не дать им подвергнуть себя чрезмерному риску. Никто не хочет этого делать "--- это неприятно, неловко и неблагодарно. Вас вряд ли отблагодарят за ваши усилия. Но альтернативы зачастую еще хуже.
		
		Существует несколько способов уберечь друзей от причинения вреда себе:
		
		\begin{itemize}[noitemsep]
			\item Организуйте безопасную поездку "--- предложите им альтернативные способы добраться до дома.
			\item Замедлите темп "--- вовлеките их в другие занятия.
			\item Удерживайте их "--- используйте любой повод, чтобы не дать им сесть на мотоцикл. Предложите еду и/или кофе, чтобы скоротать время.  Объясните, что вы беспокоитесь о том, что их могут арестовать, они могут пострадать сами или причинить вред кому-то другому.
			\item Подключайте друзей "--- используйте давление группы друзей.
		\end{itemize}
		
		Когда вы решаете вмешаться, полезно заручиться поддержкой других. Чем больше людей на вашей стороне, тем легче проявить твердость и тем сложнее мотоциклисту сопротивляться. Хотя в тот момент вам, возможно, и не поблагодарят, вам никогда не придется говорить: «Если бы я только тогда...».
\end{rus}
\end{langblock}


\clearpage
\langsection[subsection]{Fatigue}{Усталость}

\begin{langblock}{0.425}
	\begin{eng}
	
	Riding a motorcycle is more tiring than driving a car. On a long trip, you'll tire sooner than you would in a car. Avoid riding when tired.  Fatigue can affect your control of the motorcycle.
	
	\begin{itemize}[noitemsep]
		\item Protect yourself from the elements -- Wind, cold, and rain make you tire quickly. Dress warmly. A windshield is worth its cost if you plan to ride long distances.
		\item Limit your distance -- Experienced riders seldom try to ride more than about six hours a day.
		\item Take frequent rest breaks -- Stop, and get off the motorcycle at least every two hours.
		\item Don't drink or use drugs -- Artificial stimulants often result in extreme fatigue or depression when they start to wear off. Riders are unable to concentrate on the task at hand.
	\end{itemize}
	\end{eng}
	\begin{rus}
		
		Езда на мотоцикле более утомительна, чем вождение автомобиля. Во время длительной поездки вы устанете быстрее, чем в автомобиле. Старайтесь не водить мотоцикл если вы устали.  Утомленность может повлиять на вашу способность управлять мотоциклом.
		
		\begin{itemize}[noitemsep]
			\item Защищайтесь от непогоды "--- ветер, холод и дождь быстро вызывают усталость. Одевайтесь потеплее. Ветровое стекло стоит своих денег, если вы планируете ездить на дальние расстояния.
			\item Ограничивайте длительность "--- опытные мотоциклисты редко пытаются ездить более шести часов в день.
			\item Делайте частые перерывы на отдых "--- остановитесь и слезьте с мотоцикла по крайней мере раз в два часа.
			\item Не употребляйте алкоголь и наркотики "--- искусственные стимуляторы часто приводят к сильной усталости или депрессии, когда их действие начинает проходить. Теряется способность сосредоточиться на вождении.
		\end{itemize}
\end{rus}
\end{langblock}






